\chapter{Results}
\label{ch:results}

Every table in this chapter is generated from result files by the scripts named in
Appendix~\ref{app:protocol}. Where a run has not survived, the cell reads
\emph{pending} and the prose says why. A reclaimed result is not replaced with a
projection.

\section{The protocol noise floor, and the exhaustive check that is missing}
\label{sec:protocol-result}

One frozen checkpoint (large, step $10{,}000$) was scored six times under the
sampled protocol at three mask seeds and two window counts. Table~\ref{tab:protocol}
reports the six draws. They span \ProtoOldSpread{} bits/byte.

That spread is a direct measurement and is larger than the small-versus-mid
difference the fixed-data ladder was built to resolve. The distinction that
matters is between \emph{reproducibility} and \emph{accuracy}: the original
protocol was seeded, which made every number reproducible and none of them
representative. A seed selects one draw from a small sample and then reports it
with four decimal places.

The exhaustive replicate is pending. The run directory that contained runs A and
B was lost with the first sandbox, so this build cannot report the reduction that
the earlier version claimed. The table marks the two exhaustive rows
\emph{pending}; no ratio is computed and no ``sampling-free'' claim is made from a
missing run. This is a narrower result than the pre-loss manuscript, and it is the one
the surviving evidence supports.

\begin{table}[htbp]
\centering
\caption{One frozen checkpoint (large, step 10{,}000) under the sampled protocol.
The six draws span more than the largest effect the ladder was built to measure.
The exhaustive replicates are marked pending because the run directory was lost
with the first sandbox; their reduction is not claimed in this build.}
\label{tab:protocol}
\begin{tabular}{llrr}
\toprule
Protocol & Condition & Bits/byte & Scored positions \\
\midrule
\multirow{6}{*}{sampled} & seed 1234, $k=64$ & 2.3834 & 19{,}610 \\
 & seed 1234, $k=256$ & 2.3934 & 78{,}553 \\
 & seed 7, $k=64$ & 2.1682 & 19{,}555 \\
 & seed 7, $k=256$ & 2.2861 & 78{,}351 \\
 & seed 99, $k=64$ & 2.5202 & 19{,}818 \\
 & seed 99, $k=256$ & 2.3139 & 78{,}611 \\
\midrule
\multirow{2}{*}{exhaustive} & run A & \emph{pending} & \\
 & run B & \emph{pending} & \\
\bottomrule
\end{tabular}
\end{table}


\section{The fixed-data size sweep is incomplete}
\label{sec:fixed-d}

Table~\ref{tab:fixed-d} reports the fixed-budget cells that survive on this
checkout. The tiny and small rows come from sampled summaries; the mid and large
rows were reclaimed and have no held-out score on disk. The source column makes
the difference visible.

The table therefore does not establish a U-shaped scaling law. It cannot: two
sampled rows from different checkpoints are not a four-point fixed-data sweep, and
the mid and large cells that produced the turn are gone. What survives is a pair
of fixed-budget scores and a provenance fact --- the large-rung result that the
pre-loss PDF reported is not reproducible from this checkout.

This is exactly the kind of hole the thesis is required to leave open. The
temptation is to keep the old four-row curve because it is in the previous PDF;
that would make the table look complete while the underlying runs no longer exist.
The table instead names the lost rows and withholds the curve. The 4B argument in
\S\ref{sec:why-241m} does not depend on the missing large rung; it depends on the
corpus arithmetic.

\begin{table}[htbp]
\centering
\caption{Fixed-budget scores surviving on this checkout (100M tokens, 30\%
mask). The source column is the evaluation protocol that produced the score:
\emph{exhaustive} where a full-score file survived, \emph{sampled} where only a
summary survived, and \emph{lost} where the run directory was reclaimed. Sampled
and exhaustive rows are not comparable, and the mid/large cells were lost with
the first sandbox; no U-shape is claimed from this table. The generalisation gap
is held-out over train loss, both in bits/byte. Generated by
\texttt{thesis/gen\_tables.py}.}
\label{tab:fixed-d}
\begin{tabular}{lrrrrrr}
\toprule
Rung & Parameters & $D/N$ & Train & Held-out & Gap & Source \\
\midrule
tiny & 2{,}411{,}200 & 41.5 & 1.331 & \textbf{1.2937} & 1.0$\times$ & sampled \\
small & 14{,}260{,}992 & 7.0 & 1.072 & \textbf{1.0098} & 0.9$\times$ & sampled \\
mid & 33{,}696{,}768 & & & \emph{pending} & & lost \\
large & 113{,}465{,}856 & & & \emph{pending} & & lost \\
\bottomrule
\end{tabular}
\end{table}


\section{The token-budget trajectory is not available}
\label{sec:e2}

The large-rung trajectory was part of the pre-loss build. No trajectory file
survives, so Table~\ref{tab:e2} is empty. There is no statement here about whether
additional epochs help or hurt; the earlier sentence that they ``improve briefly
and then stop'' is withdrawn with the run.

\begin{table}[htbp]
\centering
\caption{The large rung as its token budget grows. The corpus holds
\CorpusTokens{} unique bytes at this build, so tokens beyond 100M are additional
epochs, not additional data.}
\label{tab:e2}
\begin{tabular}{rrr}
\toprule
Tokens seen & Epochs & Bits/byte \\
\midrule
& & \emph{pending} \\
\bottomrule
\end{tabular}
\end{table}


\section{Transfer: a matched Banking77 gain, an inconclusive AG News result}
\label{sec:transfer}

Bits per byte is a proxy. The quantity the thesis claims is decision accuracy, so
the tiny rung was fine-tuned on two classification tasks from a matched random
initialisation and from the pretrained checkpoint. Both arms use 2,000 training
rows, 500 test rows, six nominal epochs, and the same order, steps, batch,
optimiser, learning rate and training seed.

\begin{table}[htbp]
\centering
\caption{Transfer from byte-level pretraining to decision accuracy at a matched
budget: 2,000 training rows, 500 test rows, six nominal epochs, identical order,
steps, batch, optimiser, learning rate and training seed. Banking77 is the
class-balanced 77-way draw (three seeds, 500 test rows each); AG News is the
class-balanced four-way draw (three seeds, 500 unique test rows, 1,500 pooled
decisions). The AG interval is the paired crossed-seed/item bootstrap and crosses
zero; the Banking77 8,000-row arm in Table~\ref{tab:rows} is a different,
row-scaling experiment and is not this matched contrast. The per-seed AG rows are
the same six cells, shown so the pooled interval cannot hide a sign change.
Generated from \texttt{results/transfer/*.json} and the recovered six-cell summary.}
\label{tab:transfer}
\begin{tabular}{lrrr}
\toprule
Task & Random init & Pretrained & $\Delta$ accuracy \\
\midrule
banking77 & 0.322 & 0.743 & \textbf{+42.1\,pp} \\
ag-news & 0.4953 & 0.5067 & \textbf{+1.13\,pp} \\
\midrule
\multicolumn{4}{@{}l}{\emph{AG News, per seed (500 items each)}} \\
\quad ag-news s0 & 0.510 & 0.512 & +0.20\,pp \\
\quad ag-news s1 & 0.490 & 0.518 & +2.80\,pp \\
\quad ag-news s2 & 0.486 & 0.490 & +0.40\,pp \\
\bottomrule
\end{tabular}
\end{table}


On Banking77, pretraining raises accuracy by \TransferBankGain{} percentage
points: 32.2\% from random initialisation against 74.3\% from the pretrained
checkpoint, over a class-balanced draw of all seventy-seven labels. The three
pretrained seeds are 0.750, 0.732 and 0.748 against 0.350, 0.292 and 0.324; the
direction is consistent, and the gap is large relative to the seed spread.

On AG News the pooled result is \TransferAgGain{} percentage points: 50.67\% for
the pretrained arm against 49.53\% for random, over 1,500 decisions on 500 unique
test items. The paired crossed-seed/item bootstrap gives a 95\% interval of
\TransferAgLo{} to \TransferAgHi{} percentage points, so the pre-registered
verdict is \TransferAgVerdict. The per-seed rows in Table~\ref{tab:transfer}
show why: the three differences are +0.20, +2.80 and +0.40 points, all inside
their own item-bootstrap intervals, and the pooled interval is wider than the
effect it is asked to resolve. The earlier AG News estimate lay above this
interval and is withdrawn; it is not evidence for the new protocol and it is not
restated as if it were.

The 8,000-row Banking77 score arm reaches \ShipAcc\% (95\% Wilson
\ShipLo--\ShipHi) in Table~\ref{tab:rows}. That is a row-scaling result, not the
matched transfer contrast: its random-initialisation control at 8,000 rows was
not run. It is therefore not evidence that the 2,000-row gain is larger at 8k,
and it is not a single-checkpoint continuation of the six AG News transfer cells.

\section{Why the transfer probe is 2.41M, not 4B}
\label{sec:why-241m}

The 4B ceiling is a permission, not the right instrument for this corpus. The
corpus has \CorpusTokens{} tokens. At 2.41M parameters that is 677 tokens per
parameter; at 4B it is \FourBDN{} tokens per parameter, about
\FourBFactor$\times$ fewer than the 113M rung and far below the compute-optimal
line. At $D/N=20$ a 4B model would need \FourBTokensB{}B unique tokens, roughly
49 times the corpus, and about \FourBGPUHours{} GPU-hours at the measured
throughput. On a single card that is years of wall-clock before data collection.

The tiny rung is where a matched random-initialisation control is affordable and
where the model cannot store its fine-tuning set. That is the scientific reason
for the choice, not a claim that 2.41M is the product's final size. A 4B model on
1.63B tokens would be a memorisation measurement; the missing large-rung evidence
does not make that arithmetic go away.

\section{Evidence identity and the reclamation}
\label{sec:reclamation}

The first sandbox was reclaimed on 2026-10-28. The original AG News source
checkpoint \texttt{3cf421\ldots bed8} was lost with the disk and cannot be restored
byte-for-byte. The v2 corpus was rebuilt from the checked-in builder and
reproduced its recorded SHA-256 exactly. The replacement checkpoint is
\texttt{817e7f04\ldots0410}: seq 1024, batch 16, seed 0, 6,104 steps,
100,007,936 tokens, held-out 1.2937 bits/byte. It reproduces the protocol but not
the bitwise weights, because GPU execution is not bitwise-deterministic.

The six AG News cells are a valid paired comparison on the amended checkpoint;
they are not a continuation of the exact Banking77 initialisation. The Banking77
84.3\% row uses the surviving pre-loss comparator
\texttt{499fd68c\ldots ebdc}, not \texttt{817e7f04\ldots0410}. The two rows are
kept separate for that reason, and the table says so rather than averaging them
into a false single-checkpoint comparison.

Completed records were not overwritten. The recovered tree is isolated under
\texttt{results/recovered\_20261028/} and its manifest hashes the corpus, all six
AG News cells, the twenty-seven suite runs (nine tasks at three rungs) and the permutation controls. A second
reclamation now costs wall-clock, not results.

\section{The public-task suite: a capability boundary at three rungs}
\label{sec:suite}

Table~\ref{tab:suite} reports nine tasks at three rungs under one pre-registered
recipe: LR \(10^{-4}\), exactly 3{,}000 optimizer steps, batch 32, 4{,}000
stratified training rows and 2{,}000 test rows where the split allows, seed 0,
an unfrozen encoder, and each checkpoint's native context. The recipe is
identical for every cell; there is no per-task learning rate, step count, row
cap, or prompt. Every Azurite cell is measured on the dual context-candidate
path, and every cell at a given rung is scored on the identical test rows, which
is what licenses the paired comparisons quoted below. The 113.47M cell is the
fourth size in the pretraining ladder but the third measured on this suite; the
two smaller suite rungs are the 14.4M and 33.7M checkpoints.

The 113.47M column is a single new run from fresh weights, not a continuation of
the run lost on 2026-10-28: \SuiteCkptLargeTokens{} tokens, step
\SuiteCkptLargeStep, checkpoint \SuiteCkptLarge. Its cells ran in capped
concurrency waves after a nine-way wave exhausted the card's memory; the failed
attempts produced no result and are kept as an infrastructure record rather than
averaged into anything.


The final recipe is not the 500-step, \(5\times10^{-4}\) protocol that produced
the superseded table. That earlier mid-rung result was primarily an optimization
failure, not a capacity ceiling: the pre-registered four-arm sweep over learning
rate and step budget selected \(10^{-4}\) at 3{,}000 steps, and every
\(5\times10^{-4}\) arm collapsed in accuracy, NLL, and calibration. The final
table applies that single winning setting unchanged to every task. This is the
learning-rate finding, and it is why the earlier suite is superseded rather than
supplemented.

At 14.4M, \SuiteSmallAtOrAboveThree{} of \SuiteComparatorCells{} comparator
cells are not clearly below their cited point, of which \SuiteSmallClear{} are
clear wins and one is a tie. Banking77 reaches \SuiteSmallBank\%
(\SuiteSmallGapBank{} points against 78.8), Financial PhraseBank
\SuiteSmallFinancial\% (\SuiteSmallGapFinancial{} against 73.0), and DAIR
Emotion \SuiteSmallEmotion\% (\SuiteSmallGapEmotion{} against 48.0); all three
clear their points. MASSIVE intent reaches \SuiteSmallMassiveIntent\%
(\SuiteSmallGapMassiveIntent{} against Laya's \SuiteCitedMassiveIntent\% card
\citep{laya2026card}),
but its interval includes the card value, so it is a tie inside sampling noise
and not a Jev win. The four remaining comparator cells are below: CLINC150
\SuiteSmallGapClinc, AG News \SuiteSmallGapAg, XNLI \SuiteSmallGapXnli, and
Language ID \SuiteSmallGapLang{} points. MASSIVE scenario has no cited point and
is an internal diagnostic.

At 33.7M, four comparator cells are not clearly below their cited points: three
clear wins, Banking77 \SuiteMidBank\% (\SuiteMidGapBank{}), DAIR Emotion
\SuiteMidEmotion\% (\SuiteMidGapEmotion{}), and Financial PhraseBank
\SuiteMidFinancial\% (\SuiteMidGapFinancial{}), plus one tie, MASSIVE intent
\SuiteMidMassiveIntent\% (\SuiteMidGapMassiveIntent), whose interval
\SuiteMidMassiveIntentLo--\SuiteMidMassiveIntentHi{} contains Laya's
\SuiteCitedMassiveIntent\%. The remaining four are clear below: CLINC150
\SuiteMidGapClinc, AG News \SuiteMidGapAg, XNLI \SuiteMidGapXnli, and Language ID
\SuiteMidGapLang{} points. An earlier version reported three of eight here because
it counted the MASSIVE intent tie as a loss; the rule stated under the table is
the interval rule, and under it the count is four, of which three are clear. The
33.7M rung is worse than 14.4M on \SuiteLargerWorse{} of \SuiteTasks{} cells, so
this step is not a scale gain over the smaller one; \S\ref{sec:suite} shows the
113.47M rung does not restore monotonicity either.


At 113.47M, five of eight comparator cells are not clearly below their cited
points. Four are clear wins: Banking77 \SuiteLargeBank\% (\SuiteLargeGapBank),
Financial PhraseBank \SuiteLargeFinancial\% (\SuiteLargeGapFinancial), DAIR
Emotion \SuiteLargeEmotion\% (\SuiteLargeGapEmotion), and MASSIVE intent
\SuiteLargeMassiveIntent\% (\SuiteLargeGapMassiveIntent), the first rung at which
that cell clears Laya's \SuiteCitedMassiveIntent\% card rather than tying it.
CLINC150 is a tie: \SuiteLargeClinc\%, \SuiteLargeShortClinc{} points below the
cited point, with an interval of \SuiteLargeClincLo--\SuiteLargeClincHi{} that
covers the cited \SuiteCitedClinc. Three cells are clear below: AG News by
\SuiteLargeShortAg, Language ID by \SuiteLargeShortLang, and XNLI by
\SuiteLargeShortXnli{} points. MASSIVE scenario has no
cited point and remains an internal diagnostic. These counts use the
pre-registered comparator points; the strongest public Jev run on Banking77 is
above \SuiteLargeBank\% and is protocol-matched in \S\ref{sec:jev-matched},
where both pre-registered heads fall short; the best arm reaches
\QueryPctQueryTwo\%, so that cell is not a win over the best public comparator.


The path correction changes what is attributable to the model. Eight cells
changed from the candidate-only cross encoder to the dual path: AG News, DAIR
Emotion, XNLI, and Financial PhraseBank at both rungs. Their former values are
retired, not compared as capability measurements. The 14.4M exact controls are
Banking77 \SuiteSmallBank\%, CLINC150 \SuiteSmallClinc\%, and MASSIVE intent
\SuiteSmallMassiveIntent\%: the same scores and the same per-example decisions,
so those earlier measurements remain valid. The other already-dual cells move
by \SuiteSmallDeltaMassiveScenario, \SuiteSmallDeltaLang,
\SuiteMidDeltaBank, \SuiteMidDeltaClinc, \SuiteMidDeltaMassiveIntent,
\SuiteMidDeltaMassiveScenario, and \SuiteMidDeltaLang{} points. Because their
path did not change, those deltas are run-to-run variation and are not credited
to the fix.

DAIR Emotion is the largest path effect and the clearest warning against
reading a harness repair as a capability gain. The retired candidate-only
cross-encoder values were \SuiteSmallPriorEmotion\% and \SuiteMidPriorEmotion\%,
below the \SuiteEmotionMajority\% majority baseline. On the dual path the same
checkpoints and recipe score \SuiteSmallEmotion\% and \SuiteMidEmotion\%,
above the cited 48.0. The \SuiteSmallDeltaEmotion{} and \SuiteMidDeltaEmotion{}
point changes are scoring corrections, not model improvements: the pre-trained
checkpoint, the 3{,}000-step recipe, and the seed are unchanged, and only the
scoring path differs. The post-fix values are valid measurements and stand in the table, but
the jump is evidence that the earlier failure was an artifact of not reading
candidate context, not that the decider crossed a capability threshold. The
113.47M cell is \SuiteLargeVsSmallEmotion{} points from its 14.4M counterpart
(exact $p = \SuiteLargeVSSmallPEmotion$), so the correction does not create a
scale gain either.

XNLI is not the same story. The retired cross-encoder value at 14.4M was
\SuiteSmallXnliRaw\%; the flat packer lost the hypothesis entirely in
\SuiteXnliAbsentTest{} of \SuiteXnliTestN{} test pairs (and
\SuiteXnliAbsentTrain{} of \SuiteXnliTrainN{} training rows), so that number is
not a capability measurement. The dual path reaches \SuiteSmallXnli\% (95\%
Wilson \SuiteSmallXnliLo--\SuiteSmallXnliHi), a \SuiteSmallXnliGain{} point
change, but still \SuiteSmallXnliShortfall{} points below the cited 86.0\%. The
earlier field-aware-cross value of \SuiteSmallXnliPriorFix\% is also superseded
by this measurement. At 14.4M the 256-token window still truncates some pairs
and the field-aware dual packer was deliberately not attempted, so part of that
shortfall may be context. The 33.7M result removes that explanation: its
1{,}024-token context holds every test pair, and it scores \SuiteMidXnli\% (95\%
Wilson \SuiteMidXnliLo--\SuiteMidXnliHi), statistically indistinguishable from
the \SuiteXnliChance\% chance rate and \SuiteMidXnliShortfall{} points below the
cited point. XNLI is therefore a genuine model failure at both smaller
rungs, not a packing artifact at 33.7M.

XNLI at 113.47M is not a weaker version of the same failure; it is a degenerate
cell. It assigns one class to \SuiteLargeXnliTopShare\% of the test rows, its
test NLL is \SuiteLargeXnliNll{} against $\ln 3 = \SuiteLargeXnliLnK$, its mean
confidence is \SuiteLargeXnliConf{} and its ECE is \SuiteLargeXnliEce. Those are
the signatures of a constant uniform distribution: the model has learned nothing
on the task and is not wrong in an informative way. Against 14.4M it gets
\SuiteLargeOnlyXnli{} rows right that 14.4M gets wrong and \SuiteSmallOnlyXnli{}
wrong that 14.4M gets right (exact $p = \SuiteLargeVSSmallPXnli$). The smaller
rungs were not doing the task either --- their NLLs are \SuiteSmallXnliNll{} and
\SuiteMidXnliNll{} with ECE \SuiteSmallXnliEce{} and \SuiteMidXnliEce, the
profile of a late-training collapse into confident errors --- but they at least
departed from uniform. Scaling removed that departure.

The scale question the ladder was built to answer therefore has a measured
answer, and it is negative. Against 14.4M, a \SuiteLargeParamsRatio$\times$
increase in parameters at the same token budget wins \SuiteLargeWins{} cells ---
\SuiteLargeWinCells{} --- by \SuiteLargeWinMin{} to \SuiteLargeWinMax{} points,
loses \SuiteLargeLosses{} cell, XNLI, by \SuiteLargeLossAbs{} points, and leaves
\SuiteLargeUnmoved{} cells statistically unmoved. The wins sit just above the
measured run-to-run floor: across \SuiteReplicationCells{} cells executed twice
on the same checkpoint and test draw, the repeats differ by at most
\SuiteReplicationMax{} points while flipping up to \SuiteReplicationFlips{} of
2{,}000 predictions. The single loss is \SuiteLargeLossOverWinMax$\times$ the
largest win and points the wrong way. None of XNLI, AG News, CLINC150 or
Language ID is closed; CLINC150 is the only one of the four that scale moved to
a tie, and it moved it by \SuiteLargeVsSmallClinc{} points. Capacity is therefore not
sufficient for these four cells.

The checkpoint is identified by hash, not by filename: \SuiteCkptLarge, with the
step and token count above. Every cell's \texttt{evidence/provenance.json} names
the same SHA-256, and the per-example decisions, metric histories and evidence
tensors are on local disk, hash-checked against their records. The recipe and the
nine cells were fixed before the first large cell ran:
\texttt{results/suite\_large\_113m/PRE\_REGISTRATION.md} (SHA-256
\SuitePreRegSha) is in the run manifest. The nine cells cost
\SuiteLargeGpuHours{} GPU-hours of measured evaluation. Nothing in this section
was retuned after the result was seen, and no cell was dropped.


The cross-harness caveat governs every row. The cited values come from different
tokenizers, pretraining corpora, model sizes, and evaluation harnesses; the
Azurite values are single-seed point estimates with binomial intervals, not
seed-aware ones. A point above a cited number is not a demonstrated win, and an
interval overlap is not parity. MASSIVE intent's cited point is Laya's model card
\citep{laya2026card}, not a Jev number. The table is a capability boundary, not a head-to-head.

\begin{table}[htbp]
\centering
\footnotesize
\caption{9 public tasks at three rungs under one pre-registered recipe:
LR $10^{-4}$, exactly 3{,}000 optimizer steps, batch 32, 4{,}000 stratified
training rows and 2{,}000 test rows where available (Financial PhraseBank has
1{,}264 train and 1{,}000 test), seed 0, encoder unfrozen, native context of
each checkpoint (256 at 14.4M, 1{,}024 at 33.7M and 113.47M). The recipe is
identical for every cell; there is no per-task learning rate, step count, row cap
or prompt. Every Azurite cell is measured on the dual context-candidate path. The
113.47M column is a single new run from fresh weights over 1{,}134{,}658{,}560
tokens (step 69{,}255), evaluated on the same test rows as the smaller rungs.
Cited cells are published or third-party figures, not reproduced. DAIR Emotion
and XNLI at the two smaller rungs changed evaluation path in an earlier pass and
their former candidate-only values are retired. MASSIVE intent's cited point is
Laya's card, not a Jev number. No interval is shown here; 95\% Wilson intervals
and exact paired McNemar tests are in \S\ref{sec:suite} and
\texttt{results/suite\_large\_113m/COMPARISON.md}. These are cross-harness
comparisons, not a head-to-head.}
\label{tab:suite}
\setlength{\tabcolsep}{2.5pt}
\begin{tabular}{@{}lrrrrrrr@{}}
\toprule
Task & Cited & \shortstack{14.4M\\accuracy} & \shortstack{$\Delta$ vs\\cited} &
\shortstack{33.7M\\accuracy} & \shortstack{$\Delta$ vs\\cited} &
\shortstack{113.47M\\accuracy} & \shortstack{$\Delta$ vs\\cited} \\
\midrule
AG News ($K=4$) & 91.00 & 83.15 & $-7.85$ & 83.70 & $-7.30$ & 85.75 & $-5.25$ \\
DAIR Emotion ($K=6$) & 48.00 & 82.25\textsuperscript{\(\ddagger\)} & \textbf{$+34.25$} & 74.15\textsuperscript{\(\ddagger\)} & \textbf{$+26.15$} & 81.70 & \textbf{$+33.70$} \\
Banking77 ($K=77$) & 78.80 & 88.05 & \textbf{$+9.25$} & 86.50 & \textbf{$+7.70$} & 88.45 & \textbf{$+9.65$} \\
CLINC150 ($K=151$) & 89.50 & 86.80 & $-2.70$ & 82.75 & $-6.75$ & 88.50\textsuperscript{\(\dagger\)} & $-1.00$ \\
MASSIVE intent ($K=60$) & 78.30\textsuperscript{\(\ast\)} & 79.20\textsuperscript{\(\dagger\)} & $+0.90$ & 77.15\textsuperscript{\(\dagger\)} & $-1.15$ & 82.00 & \textbf{$+3.70$} \\
MASSIVE scenario ($K=18$) & --- & 84.15 & --- & 80.05 & --- & 83.90 & --- \\
XNLI ($K=3$) & 86.00 & 43.85\textsuperscript{\(\ddagger\)} & $-42.15$ & 35.00\textsuperscript{\(\ddagger\)} & $-51.00$ & 33.90 & $-52.10$ \\
Language ID ($K=20$) & 99.60 & 96.30 & $-3.30$ & 95.60 & $-4.00$ & 95.80 & $-3.80$ \\
Financial PhraseBank ($K=3$) & 73.00 & 88.10 & \textbf{$+15.10$} & 86.90 & \textbf{$+13.90$} & 90.10 & \textbf{$+17.10$} \\
\bottomrule
\end{tabular}
\par\smallskip
\footnotesize
\textsuperscript{\(\ddagger\)}Corrected dual-path cell; the text identifies its
retired cross-path value.
\quad
\textsuperscript{\(\dagger\)}95\% Wilson interval includes the cited point; not
a clear win or loss.
\quad
\textsuperscript{\(\ast\)}Laya card (78.30 for
MASSIVE intent), not a Jev number.
\quad
Clear wins against the cited point are bold; a positive $\Delta$ alone is not a win.
\end{table}


Every cell in that table is expanded in Table~\ref{tab:suite-ci}, with its Wilson
interval, its gap against the cited point and the exact paired test against
14.4M; the counts quoted above are the columns of that table, not a separate
analysis.

\begin{table}[htbp]
\centering
\footnotesize
\caption{Every 113.47M cell with its 95\% Wilson interval, its gap against the
cited point, and the exact paired McNemar test against the 14.4M cell on the same
test rows. $K$ is the number of classes, $n$ the number of scored test rows.
Verdict is the interval rule: \emph{above} and \emph{below} mean the Wilson
interval excludes the cited point, \emph{tie} that it covers it. The last two
columns compare the two Azurite rungs directly, not either against the citation.
Read from \texttt{results/suite\_large\_113m/COMPARISON.json}; the generator
recomputes every cell and refuses a mismatch.}
\label{tab:suite-ci}
\setlength{\tabcolsep}{2.5pt}
\begin{tabular}{@{}lrrrrrlrr@{}}
\toprule
Task & $K$ & $n$ & Acc. & 95\% CI & $\Delta$ cited & Verdict &
$\Delta$ vs 14.4M & $p$ \\
\midrule
AG News & 4 & 2,000 & 85.75 & 84.15--87.21 & $-5.25$ & below & $+2.60$ & $9.5\times10^{-4}$ \\
DAIR Emotion & 6 & 2,000 & 81.70 & 79.95--83.33 & $+33.70$ & above & $-0.55$ & 0.58 \\
Banking77 & 77 & 2,000 & 88.45 & 86.97--89.78 & $+9.65$ & above & $+0.40$ & 0.58 \\
CLINC150 & 151 & 2,000 & 88.50 & 87.03--89.82 & $-1.00$ & tie & $+1.70$ & 0.01 \\
MASSIVE intent & 60 & 2,000 & 82.00 & 80.26--83.62 & $+3.70$ & above & $+2.80$ & $8.4\times10^{-5}$ \\
MASSIVE scenario & 18 & 2,000 & 83.90 & 82.22--85.45 & --- & --- & $-0.25$ & 0.76 \\
XNLI & 3 & 2,000 & 33.90 & 31.86--36.00 & $-52.10$ & below & $-9.95$ & $3.6\times10^{-10}$ \\
Language ID & 20 & 2,000 & 95.80 & 94.83--96.59 & $-3.80$ & below & $-0.50$ & 0.24 \\
Financial PhraseBank & 3 & 1,000 & 90.10 & 88.09--91.80 & $+17.10$ & above & $+2.00$ & 0.05 \\
\bottomrule
\end{tabular}
\par\smallskip
\footnotesize
$\Delta$ is in percentage points. MASSIVE intent's cited point is Laya's card
(78.30), not a Jev number;
MASSIVE scenario has no published point and its gap column is empty. A tie is not
parity: it means this test draw cannot separate the two values.
\end{table}


\section{Matched-\texorpdfstring{$D/N$}{D/N} scaling is measured}
\label{sec:iso}

Table~\ref{tab:iso-dn} is the ladder the fixed-data sweep could not be: four
rungs at matched tokens per parameter, $D/N = 10$, on the v2 corpus. Tiny,
small and mid carry the exhaustive non-overlapping score; all four carry the
sampled-64 score at the identical configuration (64 windows of 1{,}024 bytes,
mask seed 1234). Held-out bits/byte falls from \IsoTinyExh{} at 2.41M parameters to
\IsoSmallExh{} at 14.26M and \IsoMidExh{} at 33.70M on the exhaustive protocol,
and from \IsoTinySamp{} to \IsoSmallSamp, \IsoMidSamp{} and \IsoLargeBpb{} on the
sampled one.

A least-squares fit of $\log \mathrm{bpb}$ on $\log N$ gives
$\alpha = \AlphaFit$ over the \AlphaPoints{} exhaustive rungs and
$\alpha = \AlphaFitSampled$ over the \AlphaPointsSampled{} sampled rungs. The
two protocols disagree on the exponent by \AlphaGap, which is small compared
with the \IsoRange-fold range of $N$ they are fitted across. At
$\alpha \approx 0.35$, a four-fold increase in parameters at fixed tokens per
parameter multiplies bits/byte by \IsoStepFactor: a \IsoStepDrop\% reduction in
held-out loss, every four-fold step, across the whole ladder.

Three limits belong on the number rather than in a footnote. The fit has four
rungs and one training seed each, so it has no interval that accounts for
run-to-run training variance; the per-point standard errors
($\approx \IsoStderrMax$ bits/byte on the exhaustive pass) are far smaller than the
between-rung differences, but they measure sampling of the held-out region, not
seeds. The large rung has only the sampled score because its exhaustive pass
was not run, so its exhaustive cell is marked absent rather than extrapolated.
And the absolute values inherit the sampled protocol's draw: all four rungs
share one draw, which is what licenses the comparison between them, but
\S\ref{sec:protocol-result} is the evidence that a different draw moves the
absolute value.

What the exponent establishes is a property of the proxy, not of the product.
Pretraining loss improves smoothly and predictably with capacity at fixed data
per parameter. Whether that improvement reaches decision accuracy is a separate
measurement, and \S\ref{sec:suite} reports it: the same three rungs, on the same
checkpoints, do not convert this loss curve into benchmark parity.

\begin{table}[htbp]
\centering
\footnotesize
\caption{Scaling at matched tokens per parameter ($D/N = 10$) on the
\CorpusTokens-token corpus. \emph{Exhaustive} is the non-overlapping full
held-out pass (9{,}967{,}360 scored bytes); \emph{sampled} is 64 windows of
1{,}024 bytes at mask seed 1234, the configuration every rung was scored with.
The large rung has only the sampled score: its exhaustive pass was not run, so
no exhaustive value is extrapolated for it. The fitted exponents are least
squares on $\log \mathrm{bpb}$ against $\log N$, over the rungs that carry each
protocol. The 256-byte context of the tiny rung bounds the sequence it was
trained on, not the number of parameters it has.}
\label{tab:iso-dn}
\begin{tabular}{lrrrrrr}
\toprule
Rung & Parameters & Context & Tokens & $D/N$ & \shortstack{Bits/byte\\exhaustive} & \shortstack{Bits/byte\\sampled-64} \\
\midrule
tiny & 2{,}411{,}200 & seq 256 & 24{,}117{,}248 & 10.00 & 1.8234 & 1.7898 \\
small & 14{,}260{,}992 & seq 1024 & 142{,}622{,}720 & 10.00 & 0.9765 & 0.9494 \\
mid & 33{,}696{,}768 & seq 1024 & 336{,}969{,}728 & 10.00 & 0.7444 & 0.7244 \\
large & 113{,}465{,}856 & seq 1024 & 1{,}134{,}673{,}920 & 10.00 & --- & 0.4515 \\
\bottomrule
\end{tabular}
\\[0.5em]
Fitted $\alpha$ in $\mathrm{bpb} = aN^{-\alpha}$: exhaustive
\textbf{0.341} over 3 rungs; sampled
\textbf{0.355} over 4 rungs.
\end{table}



\section{The context-length curve is a single point}
\label{sec:context}

Table~\ref{tab:context} has one measured row, seq 256 at 0.9901 bits/byte, and two
pending rows. The 1,024 and 4,096 runs were reclaimed. A one-point curve is not a
degradation result; the long-context claim remains a design target rather than a
measured property.

\begin{table}[htbp]
\centering
\caption{Context length at matched training tokens (14.26M parameters,
142.6M tokens each). Isolates sequence length from optimisation budget.}
\label{tab:context}
\begin{tabular}{lrr}
\toprule
Training context & Tokens & Bits/byte \\
\midrule
seq 256 & 142{,}610{,}432 & 0.9901 \\
seq 1024 & 142{,}622{,}720 & 0.9765 \\
seq 4096 & 142{,}671{,}872 & 1.2230 \\
\bottomrule
\end{tabular}
\\[0.5em]
Degradation from 256 to 4096 tokens: \textbf{+0.2329} bits/byte.
\end{table}


\section{The candidate-set sweep: width has to be trained}
\label{sec:candidates}

Every suite cell fixes its task's label width, so the suite cannot say what
happens to accuracy as the candidate set widens. That is the question the
commercial claim turns on, and \S\ref{sec:future} named a candidate-set sweep as
the direct test. The obvious implementation of that sweep is not an experiment.
With one fixed model and one fixed bank, restricting a row to a subset that still
contains the gold can only delete a wrong argmax: if the full-bank argmax is the
gold it remains the argmax of the subset, and if it is not, deleting it can only
help. Restricted accuracy is therefore at least the full-bank accuracy by
construction, whatever the model learned, and a curve produced that way would
``hold'' for a model that had learned nothing. Width has to be trained.

The sweep therefore trains a fresh decision head for each width. Every arm starts
from the same pretrained checkpoint (\SuiteCkptLarge, step \SuiteCkptLargeStep),
uses the suite's frozen fine-tuning recipe --- 4{,}000 stratified training rows,
3{,}000 optimizer steps, batch 32, LR \(10^{-4}\), AdamW, seed 0, the
checkpoint's 1{,}024-byte context --- and is tested on the suite's 2{,}000 test
rows at the width it was trained at. The only difference between arms is
\(|C|\): each training row is packed against its gold plus \(|C|-1\) distractors
drawn uniformly without replacement, with the sampling seed fixed by a hash of
task, width, step and row. Three wide-bank tasks are used: Banking77
(\(K=\EBankK\)), CLINC150 (\(K=\EClincK\)) and MASSIVE intent
(\(K=\EMassiveK\)), at \(|C| \in \{2, 8, 32, K\}\).

The design, the sampling rule and the decision rule were written and hashed
before any valid accuracy existed; both amendments, including the one that voids
the first launch, are inside the hashed file (\texttt{results/e4\_candidate\_sweep/}, SHA-256
\EPreregSha). A first launch was voided by the harness check: the width-\(K\)
Banking77 arm scored outside the 0.2-point tolerance because the decision head
was not seeded, so every arm had started from a different random
initialisation. Those runs are preserved under \texttt{\_invalid\_init} with
their reason and are excluded from every number here; the trainer now seeds
before building the head, and the design is unchanged. The pre-registered primary rule is a non-inferiority test: the
claim that accuracy holds at \(|C|=32\) survives for a task only if the paired,
one-sided 95\% bootstrap lower bound on
\(\mathrm{acc}(32)-\mathrm{acc}(K)\) is above \(-5.0\) points. The pre-registration also set a harness check: the
width-\(K\) arm, trained by the same code and schedule as the suite cell, should
reproduce the published value within 0.2 points, and any larger deviation is
reported as a discrepancy rather than repaired by substituting the published
number. Table~\ref{tab:e4} reports every arm; the full-bank arms are the same
recipe as the suite cells, not the published numbers.

\begin{table}[htbp]
\centering
\footnotesize
\caption{The candidate-set sweep at 113.47M. Each row is a separate fine-tune of
the same pretrained checkpoint, at the suite's frozen recipe (4{,}000
stratified rows, 3{,}000 steps, batch 32, LR $10^{-4}$, seed 0), differing only
in the number of candidates $|C|$ packed into every training and test row. Test
rows are the same 2{,}000 rows at every width. A fixed model restricted to a
sub-bank is not an experiment: with the gold always present, deleting candidates
can only remove wrong argmaxes, so restricted accuracy is guaranteed to be at
least the full-bank accuracy. Width is therefore trained, not simulated. Wilson
intervals are 95\%; NLL is over the sampled candidate set and is not comparable
across widths. Read from \texttt{results/e4\_candidate\_sweep/SUMMARY.json};
\texttt{scripts/e4\_summary.py} verifies every per-example record before
writing it.}
\label{tab:e4}
\setlength{\tabcolsep}{2.5pt}
\begin{tabular}{@{}lrrrrrr@{}}
\toprule
Task & $|C|$ & Accuracy & 95\% CI & NLL & $\Delta$ vs $K$ & $p$ (32 vs $K$) \\ 
\midrule
Banking77 ($K=77$) & 2 & 98.05 & 97.35--98.57 & 0.06 & $+8.90$ & --- \\
Banking77 ($K=77$) & 8 & 95.95 & 94.99--96.73 & 0.15 & $+6.80$ & --- \\
Banking77 ($K=77$) & 32 & 92.10 & 90.84--93.20 & 0.34 & $+2.95$ & $2.08\times10^{-6}$ \\
Banking77 ($K=77$) & 77 & 89.15 & 87.71--90.44 & 0.56 & --- & --- \\
CLINC150 ($K=151$) & 2 & 98.00 & 97.29--98.53 & 0.06 & $+10.30$ & --- \\
CLINC150 ($K=151$) & 8 & 95.25 & 94.23--96.10 & 0.21 & $+7.55$ & --- \\
CLINC150 ($K=151$) & 32 & 92.20 & 90.94--93.30 & 0.34 & $+4.50$ & $2.63\times10^{-12}$ \\
CLINC150 ($K=151$) & 151 & 87.70 & 86.19--89.07 & 0.60 & --- & --- \\
MASSIVE intent ($K=60$) & 2 & 96.70 & 95.82--97.40 & 0.10 & $+14.80$ & --- \\
MASSIVE intent ($K=60$) & 8 & 91.40 & 90.09--92.55 & 0.39 & $+9.50$ & --- \\
MASSIVE intent ($K=60$) & 32 & 82.85 & 81.14--84.44 & 0.91 & $+0.95$ & $2.06\times10^{-1}$ \\
MASSIVE intent ($K=60$) & 60 & 81.90 & 80.15--83.53 & 1.16 & --- & --- \\
\bottomrule
\end{tabular}
\par\smallskip
\footnotesize
$\Delta$ is in percentage points against the full-bank arm of the same task; the
$p$ column is the exact two-sided McNemar test for 32 against $K$ on the same
test rows. The full-bank arms are the harness check: the pre-registered tolerance against
the published suite cell is 0.2 points, and a larger deviation is reported as a
discrepancy rather than replaced by the published value.
\end{table}


\paragraph{What the sweep found.} On Banking77 the width-32 arm reaches
\EBankWThirtyTwo\% (95\% Wilson \EBankWThirtyTwoLo--\EBankWThirtyTwoHi) against \EBankWK\% at
width \(K\), a difference of \EBankDelta{} points with a lower bound of
\EBankLower{}; the rule's verdict is that the claim \EBankVerdict. On CLINC150
the corresponding numbers are \EClincWThirtyTwo\% against \EClincWK\%, a
difference of \EClincDelta{} points and a lower bound of \EClincLower, so the
verdict is that the claim \EClincVerdict. On MASSIVE intent the width-32 arm
reaches \EMassiveWThirtyTwo\% against \EMassiveWK\%, a difference of
\EMassiveDelta{} points (lower bound \EMassiveLower), and the verdict is that the
claim \EMassiveVerdict. Across the three tasks the pre-registered claim
\EAggregateVerdict{} on \EHolds{} of \ETasks{} tasks.

\paragraph{The curve falls with width.} The non-inferiority test passes on every
task: the 32-way head is at or above the full bank everywhere, decisively so on
Banking77 and CLINC150 and within noise on MASSIVE intent. That is not a flat curve. Accuracy declines
monotonically as the candidate set widens. On Banking77, 2 candidates reach
\EBankWTwo\%, 8 reach \EBankWEight\%, 32 reach \EBankWThirtyTwo\% and the full bank
\EBankWK\%; CLINC150 gives \EClincWTwo, \EClincWEight, \EClincWThirtyTwo{} and
\EClincWK\%, and MASSIVE intent \EMassiveWTwo, \EMassiveWEight, \EMassiveWThirtyTwo{}
and \EMassiveWK\%. The width-8 arm is ahead of the width-32 arm by
\EBankWidthEightVsThirtyTwo{} points on Banking77, \EClincWidthEightVsThirtyTwo{} on CLINC150 and
\EMassiveWidthEightVsThirtyTwo{} on MASSIVE intent.

The arms are also not interchangeable. Re-scored on the full bank, the heads
trained at width 2 fall to \EBankFullKAtTwo, \EClincFullKAtTwo{} and
\EMassiveFullKAtTwo\%, and those trained at width 32 reach \EBankFullKAtThirtyTwo,
\EClincFullKAtThirtyTwo{} and \EMassiveFullKAtThirtyTwo\% on the full bank. Training at a
width teaches a head for that width; the at-width column is the comparison the
pre-registration fixed, and a narrow head is not a full-bank model that happens
to have seen fewer candidates.

So the registered claim survives, but the stronger reading does not. The 32-way
interface is at or above the full bank on every task, but because the full bank
is the hardest point in this curve, not because width is free. Widening from 2 to 32 costs \EBankCostTwoToThirtyTwo{} points on Banking77,
\EClincCostTwoToThirtyTwo{} on CLINC150 and \EMassiveCostTwoToThirtyTwo{} on MASSIVE intent. A product
that could pose 8 candidates would score higher than one posing 32 on these
three tasks; the 32-way interface is a compromise between coverage and accuracy,
not an accuracy-neutral choice. The valuation claim needs the registered
direction, and that direction is supported; it does not get the stronger claim
that accuracy is flat as the set widens.

The harness check is reported rather than assumed. The width-\(K\) arms score
\EBankWK, \EClincWK{} and \EMassiveWK\% against published suite values of
\EBankPublished, \EClincPublished{} and \EMassivePublished\%, so the absolute
deviations are \EBankHarnessDelta, \EClincHarnessDelta{} and
\EMassiveHarnessDelta{} points. The 0.2-point tolerance was set before the run
and two arms exceed it; the largest, \EMaxHarnessDelta{} points, is still inside
the \SuiteReplicationMax-point run-to-run floor the suite measured on repeated
cells. That is a level offset of the same size as the suite's own repeat noise,
and it is reported rather than hidden by substituting the published number.
The width comparison itself is within-run --- both arms trained by the same code
on the same rows --- so a level offset against the published cell shifts the
whole curve rather than the difference the rule tests.

What the sweep establishes is bounded. It is one seed per arm, one rung, three
English tasks and 4{,}000 training rows; the per-width arms differ in
optimisation trajectory as well as width, which is inherent to training the width
rather than simulating it. It says that at 113.47M a trained 32-way interface
is at or above the full bank on these cells, and that narrower interfaces do
better still;
it does not say that accuracy is constant in width, and it does not establish
near-perfect accuracy at any width. The remaining grid --- a second seed, and the
same twelve arms at the 14.4M and 33.7M rungs --- is costed in
\S\ref{sec:future}.


\section{Summary of the measured position}

\begin{enumerate}
  \item \textbf{High confidence.} The sampled protocol has a noise floor of
  \ProtoOldSpread{} bits/byte on one frozen checkpoint, larger than the
  small-versus-mid effect the ladder was built to resolve. This is a direct
  measurement from the surviving six draws.
  \item \textbf{Moderate confidence.} At 2.41M parameters, byte-level pretraining
  transfers decision-relevant structure on 77-way Banking77:
  \TransferBankGain{} percentage points over a matched random initialisation. The
  AG News balanced result is \TransferAgGain{} percentage points with an interval
  that crosses zero; it is \TransferAgVerdict, not a gain.
  \item \textbf{Descriptive, not a win.} The 8,000-row Banking77 arm reaches
  \ShipAcc\% (95\% Wilson \ShipLo--\ShipHi), overlapping the pre-registered Jev
  interval (78.8\%, 73.1--83.6). The point estimate is above that comparator,
  the interval is not, and the rows are not shared. The protocol-matched test in
  \S\ref{sec:jev-matched} confirms the gap on shared rows: the mean-pool head
  reaches \JevMatchPctAOne\% with the same 24 examples and \JevMatchPctNoEx\%
  without them. A second, pre-registered query-anchored head recovers the
  example context to \QueryPctQueryTwo\%, still \QueryDiffAbs{} points below
  Jev; the pooling diagnosis is confirmed as a mechanism, not as a P0 win.
  \item \textbf{A capability boundary, not parity.} The nine-task suite at three
  rungs reaches, or is not clearly below, the cited point on
  \SuiteSmallAtOrAboveThree{} of \SuiteComparatorCells{} comparator cells at
  14.4M, \SuiteMidAtOrAboveThree{} at 33.7M and \SuiteLargeAtOrAbove{} at
  113.47M, of which \SuiteSmallClear, \SuiteMidClear{} and \SuiteLargeClear{} are
  clear wins. The Banking77 win is against the pre-registered 78.8\% zero-shot
  cell; a stronger public Jev run reports 92.40\%, above \SuiteLargeBank\%, so the
  flagship comparison does not survive that comparator \citep{smith2026jevbanking}.
  The protocol-matched run on the same 3,080 rows returns \JevMatchPctAOne\%
  with the mean-pool head (\JevMatchPctNoEx\% without examples); the
  query-anchored head returns \QueryPctQueryTwo\% and still misses Jev by
  \QueryDiffAbs{} points, so the gap survives both heads.
  The \SuiteLargeParamsRatio$\times$ increase from 14.4M wins
  \SuiteLargeWins{} cells by \SuiteLargeWinMin{} to \SuiteLargeWinMax{} points,
  loses XNLI by \SuiteLargeLossAbs{} points, and leaves \SuiteLargeUnmoved{}
  cells unmoved: capacity alone does not close the four cells the thesis needed. XNLI at 113.47M is a degenerate uniform predictor
  (\SuiteLargeXnliNll{} nats, one class on \SuiteLargeXnliTopShare\% of rows),
  not a weak attempt, and the \SuiteReplicationCells-cell replication puts the run-to-run floor at
  \SuiteReplicationMax{} points. DAIR Emotion's retired
  \SuiteSmallPriorEmotion\% and \SuiteMidPriorEmotion\% were a candidate-only
  scoring artifact; its dual-path values are valid measurements, but the change
  is not a capability gain.
  \item \textbf{The width test, at one rung and one seed.} On the three wide-bank
  tasks, a head trained at 32 candidates scores \EBankWThirtyTwo, \EClincWThirtyTwo{} and
  \EMassiveWThirtyTwo\% against \EBankWK, \EClincWK{} and \EMassiveWK\% at the full
  bank. Under the pre-registered 5-point non-inferiority rule the claim
  \EAggregateVerdict{} on \EHolds{} of \ETasks{} tasks. The curve falls with
  width --- 2 candidates reach \EBankWTwo, \EClincWTwo{} and \EMassiveWTwo\% --- so
  the pass is driven by 32 matching or exceeding the full bank, not by width
  being free.
  This is evidence about the 32-way interface at 113.47M, not a demonstration of
  near-perfect accuracy; the full grid is costed in \S\ref{sec:future}.
  \item \textbf{Measured scaling in the proxy.} At $D/N = 10$ the held-out loss
  follows $\mathrm{bpb} \propto N^{-\alpha}$ with $\alpha = \AlphaFit$ over the
  \AlphaPoints{} exhaustive rungs and $\alpha = \AlphaFitSampled$ over the
  \AlphaPointsSampled{} sampled rungs, across a \IsoRange-fold range of
  parameters. The proxy scales cleanly; the decision accuracy in the previous
  item does not. That contrast, not the loss curve alone, is the thesis's answer
  on scale.
  \item \textbf{Not established.} The fixed-data U-shape, the exhaustive protocol
  reduction, the context-length curve and any 4B result. These are visibly
  pending, with the reason stated in the relevant section.

\end{enumerate}
